\section{Protocol-Matched 5G-NIDD Results}

\subsection{Overall performance and family dependence}

Table~\ref{tab:familymain} reports the five-seed result for \compactmethod{}. Known-class macro-F1 is consistently high. Open-set performance, however, is family-dependent. HTTPFlood, ICMPFlood, SYNFlood, and SYNScan are well separated. UDPFlood is not: its mean AUROC is 0.4533 and its conformal recall is 0.0032. SlowrateDoS and TCPConnectScan are intermediate. Figure~\ref{fig:family} adds the seed-to-seed dispersion, and Section~\ref{sec:harddiagnostics} uses retained per-flow scores from the detection-equivalent 24-byte reference to explain the two most difficult families.

\begin{table}[H]
\caption{Five-seed 5G-NIDD performance of \compactmethod{}. Values are mean $\pm$ standard deviation across seeds. Recall and FPR use the split-conformal 5\% rule.}
\label{tab:familymain}
\centering
\small
\begin{tabular}{lrrrr}
\toprule
Held-out family & Known macro-F1 & Unknown AUROC & Unknown recall & Known FPR \\
\midrule
HTTPFlood & $0.9991\pm0.0001$ & $0.9894\pm0.0023$ & $0.9447\pm0.0123$ & $0.0504\pm0.0003$ \\
ICMPFlood & $0.9986\pm0.0001$ & $0.9966\pm0.0021$ & $1.0000\pm0.0000$ & $0.0494\pm0.0008$ \\
SlowrateDoS & $0.9991\pm0.0001$ & $0.8994\pm0.0185$ & $0.3035\pm0.1161$ & $0.0499\pm0.0006$ \\
SYNFlood & $0.9984\pm0.0003$ & $0.9924\pm0.0021$ & $1.0000\pm0.0000$ & $0.0496\pm0.0003$ \\
SYNScan & $0.9988\pm0.0003$ & $0.9958\pm0.0014$ & $1.0000\pm0.0000$ & $0.0499\pm0.0008$ \\
TCPConnectScan & $0.9988\pm0.0001$ & $0.9268\pm0.0495$ & $0.7119\pm0.2734$ & $0.0497\pm0.0002$ \\
UDPFlood & $0.9985\pm0.0001$ & $0.4533\pm0.0146$ & $0.0032\pm0.0044$ & $0.0496\pm0.0003$ \\
UDPScan & $0.9988\pm0.0001$ & $0.9643\pm0.0261$ & $0.8274\pm0.1614$ & $0.0492\pm0.0002$ \\
\midrule
All 40 trials & $0.9987\pm0.0003$ & $0.9022\pm0.1763$ & $0.7238\pm0.3714$ & $0.0497\pm0.0006$ \\
\bottomrule
\end{tabular}
\end{table}

\begin{figure}[H]
\centering
\includegraphics[width=0.86\textwidth]{figures/per_family_5g.pdf}
\caption{Five-seed open-set performance of \compactmethod{} by held-out 5G-NIDD family. Bars show means and error bars show one standard deviation. The figure exposes the family dependence hidden by the overall mean, particularly for SlowrateDoS, TCPConnectScan, UDPFlood, and UDPScan.}
\label{fig:family}
\end{figure}

\subsection{Controlled baseline comparison}

Every row of Table~\ref{tab:baselines} uses the same 40 splits and conformal 5\% rule. The centralized Random Forest is a strong full-feature reference. Its mean recall is higher than that of \compactmethod{}, whereas its mean AUROC is lower; neither open-set difference is significant after Holm correction. \compactmethod{} is therefore not claimed to dominate the centralized model. Its principal advantage is a 16-byte evidence payload instead of a 352-byte transformed vector, a 95.5\% reduction.

The all-source logit average performs poorly on known classes under the source-aware non-IID ownership. This result demonstrates that simple averaging is not a suitable distributed comparator in this protocol, rather than implying that all distributed IDS methods fail. FedAvg preserves moderate open-set ranking but loses known-class macro-F1.

\begin{table}[H]
\caption{Protocol-matched comparison over eight 5G-NIDD holdouts and five seeds ($n=40$). All recall and FPR values use the split-conformal 5\% rule. ``Bytes'' denotes inference-time evidence/feature payload; FedAvg exchanges model updates during training and is therefore not assigned a per-flow payload.}
\label{tab:baselines}
\centering
\scriptsize
\begin{tabular}{L{0.25\textwidth}rrrrrr}
\toprule
Method & Bytes & Known F1 & AUROC & Recall & Mean FPR & Coordinator $\mu$s/flow \\
\midrule
\compactmethod{} & 16 & 0.9987 & 0.9022 & 0.7238 & 0.0497 & 0.959 \\
\referencemethod{} & 24 & 0.9987 & 0.9022 & 0.7238 & 0.0497 & 0.929 \\
\dhmethod{} & 24 & 0.9987 & 0.9050 & 0.7042 & 0.0498 & 0.929 \\
All-source logit average & 36 & 0.0753 & 0.8505 & 0.5104 & 0.0499 & 0.355 \\
FedAvg linear & -- & 0.9013 & 0.9004 & 0.6600 & 0.0500 & 0.232 \\
Central Random Forest & 352 & 0.9988 & 0.8873 & 0.7986 & 0.0348 & 0.339 \\
Central Extra Trees & 352 & 0.9941 & 0.8868 & 0.6182 & 0.0489 & 0.466 \\
Central RF + EVT tail & 352 & 0.9988 & 0.7886 & 0.3979 & 0.0503 & 0.339 \\
\bottomrule
\end{tabular}
\end{table}

\subsection{Communication--performance frontier}

The 16-byte encoding retains the performance of the 24-byte reference. Across 40 matched trials, its mean differences relative to 24B are $-2.5\times10^{-7}$ macro-F1, $6.7\times10^{-6}$ AUROC, and $2.4\times10^{-5}$ recall. The corresponding paired $t$-test p-values are 0.323, 0.161, and 0.600; the Wilcoxon p-values are 0.317, 0.646, and 0.730. The quantized encoding is therefore the compact Pareto operating point in the evaluated family. The 20- and 24-byte rows carry the same predictive payload and differ only in framing; their quality is accordingly identical.

\begin{table}[H]
\caption{Five-seed communication ablation over 40 matched trials. The full-feature row transmits all transformed float32 features.}
\label{tab:pareto}
\centering
\begin{tabular}{lrrrr}
\toprule
Encoding & Bytes & Known F1 & Unknown AUROC & Unknown recall \\
\midrule
Class+anomaly & 12 & 0.9987 & 0.9007 & 0.7117 \\
\compactmethod{} & 16 & 0.9987 & 0.9022 & 0.7238 \\
Class+proxy & 18 & 0.9988 & 0.8980 & 0.7004 \\
X-MAG-COS-20B & 20 & 0.9987 & 0.9022 & 0.7238 \\
\referencemethod{} & 24 & 0.9987 & 0.9022 & 0.7238 \\
X-MAG-COS-30B & 30 & 0.9988 & 0.8991 & 0.6883 \\
Central full features & 352 & 0.9988 & 0.8873 & 0.7986 \\
\bottomrule
\end{tabular}
\end{table}

\begin{figure}[H]
\centering
\includegraphics[width=0.82\textwidth]{figures/pareto_auroc.pdf}
\caption{Communication--performance frontier over 40 matched trials. The horizontal axis uses a logarithmic scale to display the 12--36-byte evidence messages together with the 352-byte centralized feature vector. The 16-byte message is the smallest evaluated full-content encoding and is indistinguishable in quality from the 24-byte reference.}
\label{fig:pareto}
\end{figure}

\subsection{Hard-family diagnostic evidence}
\label{sec:harddiagnostics}

The communication frontier establishes that quantization is not responsible for the difficult-family behavior: the 16-byte and 24-byte full-content encodings have indistinguishable aggregate detection quality. The retained per-flow diagnostic arrays were generated for the 24-byte float32 reference and are therefore used only to examine the failure mechanism. Receiver operating characteristic (ROC) curves are threshold-free. For the score-distribution and missed-decision analyses, a validation 95th-percentile threshold was used; its mean known-traffic false-positive rates were 0.0496 for UDPFlood and 0.0500 for SlowrateDoS, closely matching the nominal 5\% operating point.

Figure~\ref{fig:hardroc} shows qualitatively different behavior for the two families. SlowrateDoS retains useful ranking information (mean AUROC $0.8994\pm0.0185$), although its thresholded recall is only $0.3035\pm0.1162$. UDPFlood is below random ranking on average (AUROC $0.4533\pm0.0146$) and is almost never rejected (recall $0.0032\pm0.0044$). Figure~\ref{fig:hardhist} explains this contrast: SlowrateDoS scores are shifted toward the upper part of the range but still overlap known traffic near the operating threshold, whereas UDPFlood scores overlap the low-score known-traffic region extensively.

\begin{figure}[H]
\centering
\subfloat[UDPFlood.]{\includegraphics[width=0.49\textwidth]{figures/roc_udpflood.pdf}}
\hfill
\subfloat[SlowrateDoS.]{\includegraphics[width=0.49\textwidth]{figures/roc_slowratedos.pdf}}
\caption{Five-seed open-set ROC diagnostics for the two difficult 5G-NIDD holdouts using the detection-equivalent 24-byte float32 reference. The solid curve is the seed mean, the band is one standard deviation, and the diagonal denotes random ranking. The vertical line marks a false-positive rate of 0.05 for visual reference.}
\label{fig:hardroc}
\end{figure}

\begin{figure}[H]
\centering
\subfloat[UDPFlood.]{\includegraphics[width=0.49\textwidth]{figures/score_histogram_udpflood.pdf}}
\hfill
\subfloat[SlowrateDoS.]{\includegraphics[width=0.49\textwidth]{figures/score_histogram_slowratedos.pdf}}
\caption{Composite-score distributions for known test traffic and the held-out families, pooled from five seed evaluations of the 24-byte reference. UDPFlood occupies the same low-score region as much of the known traffic, whereas SlowrateDoS is more strongly shifted but retains substantial overlap around the operating threshold.}
\label{fig:hardhist}
\end{figure}

The class assignments of missed held-out decisions further distinguish the mechanisms (Figure~\ref{fig:hardassign}). Across the five seed evaluations, all 2,279,445 accepted UDPFlood decisions were assigned to Benign. Among 254,662 accepted SlowrateDoS decisions, 253,875 (99.691\%) were assigned to HTTPFlood and 787 (0.309\%) to Benign. These are repeated seed-trial decisions over the same held-out records rather than counts of unique flows. The SlowrateDoS error therefore supports absorption into a semantically related application-layer flooding class. The UDPFlood error does not: it is a joint representation and rejection failure in which the held-out traffic appears benign to the compact known-class head and receives low composite scores.

\begin{figure}[H]
\centering
\includegraphics[width=0.78\textwidth]{figures/accepted_unknown_assignments.pdf}
\caption{Known-class assignments for held-out decisions that fell below the validation 95th-percentile rejection threshold, aggregated over five seeds. All missed UDPFlood decisions were assigned to Benign, whereas 99.691\% of missed SlowrateDoS decisions were assigned to HTTPFlood.}
\label{fig:hardassign}
\end{figure}
